% GENERATED FILE. Edit canonical lessons, then run npm run generate:books.
% canonical-source-hash: fnv1a64:cd1e6b47849c762e
% canonical-lessons: HI-C417-disa, HI-C417-taraf, HI-C417-or, HI-C417-marg, HI-C417-rukh

\chapter{Which Way?}
\label{ch:hi-which-way}
\hlchaptermodality{417}

\begin{chapteropening}
\textbf{By the end of this chapter:} \emph{I can ask which direction and which route to take, and say which way the wind is blowing.}

Complete the last lesson of chapter 417: I can ask which direction and which route to take, and say which way the wind is blowing.
\end{chapteropening}

\section[\texorpdfstring{\hi{दिशा} (diśā)}{दिशा (diśā)}]{\hi{दिशा} (diśā) — a direction}

\label{lesson:HI-C417-disa}



\begin{quote}
Before the new one: say the Hindi for to spend (time), then the Hindi for to pass (time).
\end{quote}

\subsection*{You'll want to know: \hi{दिशा}}
\textbf{\hi{दिशा}} — \emph{diśā} — \textquotedblleft{}a direction\textquotedblright{}.

\begin{quote}
\hi{स्टेशन} \hi{किस} \hi{दिशा} \hi{में} \hi{है}?
\end{quote}

\begin{quote}
\emph{sṭeśan kis diśā meṁ hai?}
\end{quote}

\begin{quote}
Which direction is the station in?
\end{quote}

Feminine.

\begin{cousinweb}
\textbf{\hi{दिशा}} (\emph{diśā}) is a tatsama, Sanskrit \emph{diśā}, \textquotedblleft{}a direction, a quarter of the sky\textquotedblright{} (CDIAL 6339); the same entry lists Hindi \emph{dis}, the form worn down in speech.
\end{cousinweb}

\begin{grammarlens}[title={What you've built}]
One more word for getting around.
\end{grammarlens}

\subsection*{Your turn}
\begin{itemize}
\raggedright
  \item \emph{Say it:} \emph{diśā}
  \item \emph{Say it:} \emph{diśā}, once more
  \item \emph{Say it:} \emph{sṭeśan kis diśā meṁ hai?} — Which direction is the station in?
  \item \emph{Recall:} say the Hindi for to spend (time), then the Hindi for to pass (time), then say \emph{diśā} again
\end{itemize}

\begin{tcolorbox}[breakable,colback=teal!4,colframe=teal!35!black,title={Before you move on}]
What does \hi{दिशा} mean? (\textquotedblleft{}A direction\textquotedblright{}.) Say it once more.
\end{tcolorbox}
\section[\texorpdfstring{\hi{तरफ़} (taraf)}{तरफ़ (taraf)}]{\hi{तरफ़} (taraf) — a side (of a place); towards (with \hi{की})}

\label{lesson:HI-C417-taraf}



\begin{quote}
Before the new one: say the Hindi for to pass (time), then the Hindi for a direction.
\end{quote}

\subsection*{You'll want to know: \hi{तरफ़}}
\textbf{\hi{तरफ़}} — \emph{taraf} — \textquotedblleft{}a side (of a place); towards (with \hi{की})\textquotedblright{}.

\begin{quote}
\hi{नदी} \hi{की} \hi{दूसरी} \hi{तरफ़} \hi{एक} \hi{मंदिर} \hi{है।}
\end{quote}

\begin{quote}
\emph{nadī kī dūsrī taraf ek mandir hai.}
\end{quote}

\begin{quote}
There is a temple on the other side of the river.
\end{quote}

Feminine.

\hi{की} \hi{तरफ़}, towards: \hi{स्टेशन} \hi{की} \hi{तरफ़} \hi{चलिए}, walk towards the station.

\begin{cousinweb}
\textbf{\hi{तरफ़}} (\emph{taraf}) is usually traced to Arabic \emph{ṭaraf}, \textquotedblleft{}a side, an edge\textquotedblright{}, which reached Hindi through Persian.
\end{cousinweb}

\begin{grammarlens}[title={What you've built}]
One more word for getting around.
\end{grammarlens}

\subsection*{Your turn}
\begin{itemize}
\raggedright
  \item \emph{Say it:} \emph{taraf}
  \item \emph{Say it:} \emph{taraf}, once more
  \item \emph{Say it:} \emph{nadī kī dūsrī taraf ek mandir hai.} — There is a temple on the other side of the river.
  \item \emph{Recall:} say the Hindi for to pass (time), then the Hindi for a direction, then say \emph{taraf} again
\end{itemize}

\begin{tcolorbox}[breakable,colback=teal!4,colframe=teal!35!black,title={Before you move on}]
What does \hi{तरफ़} mean? (\textquotedblleft{}A side (of a place); towards (with \hi{की})\textquotedblright{}.) Say it once more.
\end{tcolorbox}
\section[\texorpdfstring{\hi{ओर} (or)}{ओर (or)}]{\hi{ओर} (or) — a side; towards (with \hi{की})}

\label{lesson:HI-C417-or}



\begin{quote}
Before the new one: say the Hindi for a direction, then the Hindi for a side (of a place).
\end{quote}

\subsection*{You'll want to know: \hi{ओर}}
\textbf{\hi{ओर}} — \emph{or} — \textquotedblleft{}a side; towards (with \hi{की})\textquotedblright{}.

\begin{quote}
\hi{हम} \hi{पहाड़} \hi{की} \hi{ओर} \hi{चल} \hi{पड़े।}
\end{quote}

\begin{quote}
\emph{ham pahāṛ kī or cal paṛe.}
\end{quote}

\begin{quote}
We set off towards the mountain.
\end{quote}

Feminine.

\begin{cousinweb}
\textbf{\hi{ओर}} (\emph{or}) is a tadbhava of Sanskrit \emph{avarā}, \textquotedblleft{}the near side\textquotedblright{} (CDIAL 812), which lists Hindi \emph{or}, \textquotedblleft{}side, nearness\textquotedblright{}. It is the inherited partner of \hi{तरफ़}: \hi{की} \hi{ओर}, towards.
\end{cousinweb}

\begin{grammarlens}[title={What you've built}]
One more word for getting around.
\end{grammarlens}

\subsection*{Your turn}
\begin{itemize}
\raggedright
  \item \emph{Say it:} \emph{or}
  \item \emph{Say it:} \emph{or}, once more
  \item \emph{Say it:} \emph{ham pahāṛ kī or cal paṛe.} — We set off towards the mountain.
  \item \emph{Recall:} say the Hindi for a direction, then the Hindi for a side (of a place), then say \emph{or} again
\end{itemize}

\begin{tcolorbox}[breakable,colback=teal!4,colframe=teal!35!black,title={Before you move on}]
What does \hi{ओर} mean? (\textquotedblleft{}A side; towards (with \hi{की})\textquotedblright{}.) Say it once more.
\end{tcolorbox}
\section[\texorpdfstring{\hi{मार्ग} (mārg)}{मार्ग (mārg)}]{\hi{मार्ग} (mārg) — a route; a way (formal)}

\label{lesson:HI-C417-marg}



\begin{quote}
Before the new one: say the Hindi for a side (of a place), then the Hindi for towards.
\end{quote}

\subsection*{You'll want to know: \hi{मार्ग}}
\textbf{\hi{मार्ग}} — \emph{mārg} — \textquotedblleft{}a route; a way (formal)\textquotedblright{}.

\begin{quote}
\hi{यह} \hi{बस} \hi{किस} \hi{मार्ग} \hi{से} \hi{जाती} \hi{है}?
\end{quote}

\begin{quote}
\emph{yah bas kis mārg se jātī hai?}
\end{quote}

\begin{quote}
Which route does this bus take?
\end{quote}

Masculine.

\begin{cousinweb}
\textbf{\hi{मार्ग}} (\emph{mārg}) is a tatsama, Sanskrit \emph{mārga}, first the track of a wild animal and then any road or way (Monier-Williams); street signs and bus boards use it, and it is the second half of \hi{राजमार्ग}.
\end{cousinweb}

\begin{grammarlens}[title={What you've built}]
One more word for getting around.
\end{grammarlens}

\subsection*{Your turn}
\begin{itemize}
\raggedright
  \item \emph{Say it:} \emph{mārg}
  \item \emph{Say it:} \emph{mārg}, once more
  \item \emph{Say it:} \emph{yah bas kis mārg se jātī hai?} — Which route does this bus take?
  \item \emph{Recall:} say the Hindi for a side (of a place), then the Hindi for towards, then say \emph{mārg} again
\end{itemize}

\begin{tcolorbox}[breakable,colback=teal!4,colframe=teal!35!black,title={Before you move on}]
What does \hi{मार्ग} mean? (\textquotedblleft{}A route; a way (formal)\textquotedblright{}.) Say it once more.
\end{tcolorbox}
\section[\texorpdfstring{\hi{रुख़} (rukh)}{रुख़ (rukh)}]{\hi{रुख़} (rukh) — a direction; the way something faces}

\label{lesson:HI-C417-rukh}



\begin{quote}
Before the new one: say the Hindi for towards, then the Hindi for a route.
\end{quote}

\subsection*{You'll want to know: \hi{रुख़}}
\textbf{\hi{रुख़}} — \emph{rukh} — \textquotedblleft{}a direction; the way something faces\textquotedblright{}.

\begin{quote}
\hi{शाम} \hi{को} \hi{हवा} \hi{का} \hi{रुख़} \hi{बदल} \hi{गया।}
\end{quote}

\begin{quote}
\emph{śām ko havā kā rukh badal gayā.}
\end{quote}

\begin{quote}
In the evening the wind changed direction.
\end{quote}

Masculine.

\begin{cousinweb}
\textbf{\hi{रुख़}} (\emph{rukh}) is usually traced to Persian \emph{rukh}, \textquotedblleft{}the face, the cheek\textquotedblright{}: the way a face is turned, and so the direction of the wind, a road or a person's mood.
\end{cousinweb}

\begin{grammarlens}[title={What you've built}]
One more word for getting around.
\end{grammarlens}

\subsection*{Your turn}
\begin{itemize}
\raggedright
  \item \emph{Say it:} \emph{rukh}
  \item \emph{Say it:} \emph{rukh}, once more
  \item \emph{Say it:} \emph{śām ko havā kā rukh badal gayā.} — In the evening the wind changed direction.
  \item \emph{Recall:} say the Hindi for towards, then the Hindi for a route, then say \emph{rukh} again
\end{itemize}

\begin{tcolorbox}[breakable,colback=teal!4,colframe=teal!35!black,title={Before you move on}]
What does \hi{रुख़} mean? (\textquotedblleft{}A direction; the way something faces\textquotedblright{}.) Say it once more.
\end{tcolorbox}
